\subsection{System architecture}
\label{sec:architecture}

VIGÍA is organised in three tiers. Detectors are interchangeable; the
validator is shared and hazard-agnostic. Adding a hazard requires a detector
class and a configuration entry, and no change to the validator.

\begin{figure}[htbp]
\centering
\begin{tikzpicture}[
  font=\small,
  box/.style={draw=rulegrey, fill=tierfill, rounded corners=1pt,
              minimum height=7mm, inner xsep=6pt, align=center},
  tier/.style={draw=none, font=\scriptsize\scshape, text=notegrey},
  arrow/.style={-{Latex[length=4pt]}, draw=notegrey},
]
\node[tier] (t0) at (0,3.3) {Tier 0 — context gate};
\node[tier] (t1) at (5.1,3.3) {Tier 1 — detectors};
\node[tier] (t2) at (10.2,3.3) {Tier 2 — validator};

\node[box, text width=25mm] (gate) at (0,1.4)
  {camera registry\\[1pt]\scriptsize declares context,\\[-1pt]\scriptsize
   viewpoint, hazards};

\node[box, text width=27mm] (d1) at (5.1,2.6) {fire \& smoke};
\node[box, text width=27mm] (d2) at (5.1,1.9) {flood / water level};
\node[box, text width=27mm] (d3) at (5.1,1.2) {people in water};
\node[box, text width=27mm] (d4) at (5.1,0.5) {building access};
\node[box, text width=27mm] (d5) at (5.1,-0.2) {road incidents};

\node[box, text width=27mm] (val) at (10.2,1.4)
  {confidence\\[-1pt] persistence (IoU)\\[-1pt] location cooldown};

\node[align=left, font=\scriptsize, text=notegrey] (out) at (10.2,-0.35)
  {\textsc{confirmed} $\rightarrow$ alert\\ \textsc{rejected} $\rightarrow$
   logged with the level\\ that rejected it};

\foreach \d in {d1,d2,d3,d4,d5} { \draw[arrow] (gate) -- (\d); }
\foreach \d in {d1,d2,d3,d4,d5} { \draw[arrow] (\d) -- (val); }
\draw[arrow] (val) -- (out);

\node[align=left, font=\scriptsize, text=notegrey] at (0,-0.1)
  {forest, street, river,\\ coast, motorway,\\ earthquake zone, \dots};
\draw[arrow] (0,0.5) -- (gate);
\end{tikzpicture}
\caption{The three tiers. The gate decides which detectors run for a given
camera; the detectors run independently and in parallel; one shared validator,
configured per hazard but never specialised per hazard, decides which
detections become events.}
\label{fig:architecture}
\end{figure}

\subsubsection{Tier 0 — the context gate}
\label{sec:gate}

Each camera is registered once with a declared context. The gate looks up
which detectors apply to that context and runs only those.

\begin{definition}[Gate]
\label{def:gate}
Let $\contexts$ be the set of camera contexts, $\hazards$ the set of hazards,
$\cameras$ the registered cameras, and
\[
  \viewpoints = \{\,\textsc{ground},\ \textsc{oblique},\ \textsc{nadir}\,\}
\]
the set of viewpoints. Each camera $c$ carries
a context $\chi(c) \in \contexts$ and a viewpoint $\nu(c) \in \viewpoints$.
The gate is a total map
\[
  G : \contexts \longrightarrow 2^{\hazards},
\]
and the detectors constructed for camera $c$ are
$\detectors(c) = \{\, d_h : h \in G(\chi(c)) \,\}$.
\end{definition}

The gate is declarative, not learned. We deliberately do not classify a frame
and route it: that introduces a misrouting failure mode on the critical path
and saves little. Every operational deployment we surveyed works this way — a
ridgeline wildfire camera never runs drowning detection. Published work on
cascaded gating in video pipelines~\citep{kang2017noscope,mao2018catdet}
reports 5--13$\times$ compute reduction at minimal accuracy cost; here the
gate is what makes five hazards affordable on one machine.

The gate's entire intelligence is \cref{tab:gate}, which is small enough to
audit at a glance. That is the trade being made: a wrong entry here is visible
in a YAML file and fixable by a human, whereas a misrouting classifier fails
silently.

\begin{table}[htbp]
\centering\small
\caption{The context-to-hazard map. The bias is toward including a hazard when
it is plausible, because a missed hazard is a life-safety failure whereas an
unnecessary detector costs only the compute the gate was saving. Where that
bias was resisted — motorway, which excludes flood — the reason is that
claiming a hazard outside a detector's measured envelope is worse than not
claiming it.}
\label{tab:gate}
\begin{tabularx}{\textwidth}{l l L}
\toprule
Context & Hazards enabled & Reasoning \\
\midrule
forest             & fire                              & PyroNear's stated envelope: fixed-mount wide-field, daylight \\
wildland interface & fire, road                        & Settlement meets vegetation, plus the roads people flee along \\
street             & flood, road, building access      & The DANA case: water rising, vehicles caught, people trapped \\
urban block        & flood, building access            & Dense built environment \\
river              & flood, drowning                   & Level and rate of rise, plus anyone swept in \\
coast              & drowning                          & Open water; the SeaDronesSee envelope \\
motorway           & road                              & Deliberately narrow: flood is plausible but outside the flood detector's calibrated envelope \\
earthquake zone    & building access                   & The DRespNeT envelope \\
UAV search         & building access, drowning         & Aerial search and rescue; both detectors are UAV-trained \\
\bottomrule
\end{tabularx}
\end{table}

\paragraph{The saving is reported, not assumed.}
For a registry of $\cameras$ cameras and $|\hazards|$ available hazards, the
gate's reduction in detector invocations per frame is
\begin{equation}
  \rho \;=\; 1 - \frac{\sum_{c \in \cameras} |G(\chi(c))|}
                      {|\cameras|\,|\hazards|}.
  \label{eq:gate-saving}
\end{equation}
On the demonstration registry shipped with the system (\GateCameras{}
cameras and $|\hazards| = \GateHazards$ available hazards), the gate requests
\GateGated{} detector invocations per frame instead of \GateUngated{}, so
$\rho = \GateReduction\,\%$, and only \GateModelsLoaded{} of the
\GateHazards{} detector types are ever constructed. The unrequested models are
never constructed at all, so the saving is in memory as well as compute: a
deployment of only ridgeline cameras loads one ONNX session, not five.

\subsubsection{Viewpoint envelopes}
\label{sec:viewpoint}

The gate also declares \emph{viewpoint}, because a detector's training
viewpoint is part of its operating envelope and nothing else in the system
captured it.

\begin{definition}[Admissibility]
\label{def:admissible}
Each detector $d$ declares an envelope $V(d) \subseteq \viewpoints$. A camera
$c$ may be paired with $d$ only if
\[
  \mathrm{adm}(c, d) \iff \nu(c) \in V(d),
\]
and $\mathrm{CameraPipeline}(c)$ raises unless
$\mathrm{adm}(c,d)$ holds for every $d \in \detectors(c)$.
\end{definition}

This was added after a measured failure, not as a precaution.

\begin{recorded}{A confident wrong number}
The flood segmenter is trained on ATLANTIS~\citep{erfani2022atlantis}, which
is ground-level and oblique photography of water bodies. Evaluated on
straight-down UAV imagery it reported up to $0.397$ water coverage on frames
whose actual water was a narrow canal, tinting large areas of mown grass as
water. It did not fail loudly. It returned a confident wrong number, which for
a life-safety component is the worst available failure. Viewpoint is therefore
declared on every camera and every detector, and the pipeline refuses the
pairing before any model is loaded — the same reasoning as the licence and
privacy guards: make the mistake impossible instead of documented. A test
fails the build if any detector stops declaring an envelope, so a sixth hazard
cannot be added without stating where it works.
\end{recorded}

\begin{recorded}{Failing loudly}
An unregistered camera raises instead of returning an empty hazard set.
Returning an empty set would turn a typo into a camera that watches for
nothing and reports success — the worst available failure mode for this
system. Enforced by \texttt{tests/test\_registry.py}.
\end{recorded}

\subsubsection{Tier 1 — specialist detectors}

Independent models running in parallel, rather than one shared-backbone
multi-task model. A unified model would require a single coherent label space
and jointly annotated data; these hazards come from separate datasets with
incompatible labels. It would also couple training, so that retraining fire
could silently degrade flood.

\subsubsection{Licensing architecture}
\label{sec:licensing}

VIGÍA ships under Apache-2.0. Several source models were fine-tuned using
Ultralytics tooling, which is licensed AGPL-3.0 — a licence whose terms would
extend to this entire project were the package imported at runtime.

The resolution is to run every model through ONNX Runtime and never import the
\texttt{ultralytics} package in the runtime path. PyroNear's own edge engine
does exactly this, which is how it remains Apache-2.0 while running a YOLO
architecture. Any \texttt{.pt}~$\rightarrow$~\texttt{.onnx} conversion is an
offline build step confined to \texttt{tools/}, executed in a separate virtual
environment.

This is enforced automatically, not by discipline:
\texttt{tools/check\_licence.py} parses every file under \texttt{vigia/} with
Python's AST module and fails the build if a forbidden package is imported. It
parses rather than greps, so a comment mentioning the package does not trip
it.

\subsubsection{Privacy by design}
\label{sec:privacy}

The EU AI Act's high-risk provisions~\citep{euaiact} took full effect in
August 2026. Annex III classifies biometric identification systems as
high-risk, and Article 5(1)(h) prohibits real-time remote biometric
identification in publicly accessible spaces for law enforcement, with narrow
exceptions.

VIGÍA identifies nobody, and this is structural, not configured.

\begin{table}[htbp]
\centering\small
\caption{Four privacy commitments and the mechanism that enforces each. A
commitment that lives only in a document is a configuration flag waiting to
happen.}
\label{tab:privacy}
\begin{tabularx}{\textwidth}{p{4.1cm} L}
\toprule
Commitment & How it is enforced \\
\midrule
No biometric code path &
\texttt{tools/check\_privacy.py} parses every file under \texttt{vigia/},
\texttt{tools/}, \texttt{scripts/} and \texttt{eval/} and fails on any import
of a face or re-identification library, any use of OpenCV's face APIs, or any
function or class whose name describes identity work. Unlike the licence
guard, \texttt{tools/} is not exempt: an offline script that built a face
index would breach the commitment just as surely. \\
\addlinespace
Only the event and one evidence frame leave &
The \texttt{Alert} type has no field for a clip and no writer for one. A test
asserts the absence, so adding one becomes a visible change to a documented
boundary rather than a quiet tweak. \\
\addlinespace
No raw video retention &
\texttt{RetentionPolicy} raises on construction if raw video storage is
requested. The rolling buffer is a fixed-length in-memory deque sized from the
camera's declared retention window, and no code path writes video to disk. \\
\addlinespace
Presence and location, never identity &
Person-shaped detections are blurred inside their own boxes before any
evidence frame is written, on a copy so the source frame is never mutated. The
box stays visible: an operator needs to know that someone is there and where,
never who. \\
\bottomrule
\end{tabularx}
\end{table}

The redaction list is keyed on detector class names, and a test fails if a
detector declares a person-depicting class that is not on it — so a sixth
hazard cannot introduce an unredacted person class by omission.

\begin{recorded}{Guards are tested against violations}
A guard that cannot fail is worthless. The privacy guard was run against
deliberately planted violations, which is how two gaps were found: class names
such as \texttt{PersonReidentifier} and \texttt{GaitSignatureExtractor}
initially escaped, because the patterns matched whole words containing
underscores. Matching now strips underscores before comparison. The same
reasoning applies to the ONNX parity check of \cref{sec:backends}.
\end{recorded}

\subsubsection{Runtime pipeline}
\label{sec:pipeline}

The three tiers are joined by \texttt{vigia/pipeline.py}: capture, gate,
detect in parallel, validate, emit. It decides none of those things itself —
the registry decides which detectors run, the validator decides what is an
alert, and the egress module decides what may leave. Keeping those three out
of the pipeline is what allows each to be tested and reported separately, and
it is why adding a sixth hazard touches none of this code.

Concurrency follows the pattern reported for a comparable
system~\citep{multitasksystem}: one
thread per active detector over bounded queues, with a single fusion stage.
The detectors are independent ONNX Runtime sessions, which release the GIL
during inference — where essentially all of the time goes — so they run
genuinely in parallel. Across cameras the detector objects are shared while
the validators are emphatically not: validator state is per camera per hazard,
and sharing it would let one camera's tracks confirm another camera's events.

The queues are bounded by design. An unbounded queue converts a detector that
cannot keep up into unbounded memory growth and ever-increasing latency, which
in a life-safety system means alerting about something that stopped being true
minutes ago. Bounded queues convert the same condition into dropped frames
that are counted and reported. How the backpressure resolves depends on the
source, and the distinction matters for the integrity of every figure here:

\begin{itemize}
  \item \textbf{Live sources never block.} A detector that falls behind costs
        frames, not latency, and stale frames are discarded in favour of the
        newest available.
  \item \textbf{File sources always block.} An evaluation replay or a curated
        demonstration clip must be deterministic; dropping frames would make a
        published measurement depend on how busy the machine happened to be.
\end{itemize}

\begin{recorded}{Found by measurement, not foresight}
With non-blocking dispatch on both source types, a 16-frame clip through a
\SI{193}{\milli\second} detector processed 3 frames and silently dropped 13. A
replay that quietly skips 80\,\% of its input would have produced plausible
and meaningless numbers.
\end{recorded}

On the reference machine a single ridgeline camera replaying a 40-frame
sequence through the fire detector sustains 8--14 frames per second end to end
at \SI{117}{\milli\second} median detector latency, processing every frame
with none dropped.
